% !TeX root = ../main.tex
%%% ---------------------------------------------------------------------------
%%% 实验结果
%%%
%%% 幻灯片上的表格要比论文里的**再删一半**：只留支撑当前论点的行与列，
%%% 完整表格放备用页。数字用 \alert 或加粗指出该看哪一个。
%%% ---------------------------------------------------------------------------
\section{实验结果}

\begin{frame}{mAP 提升 3.7 个百分点，小目标提升 6.2 个百分点}
  \begin{table}
    \centering
    \begin{tabular}{l S[table-format=2.1] S[table-format=2.1] S[table-format=2.0]}
      \toprule
      {方法} & {mAP / \%} & {AP\textsubscript{S} / \%} & {速度 / FPS} \\
      \midrule
      RetinaNet    & 69.7 & 43.0 & 28 \\
      FPN + ATSS   & 74.7 & 48.5 & 33 \\
      \midrule
      \textbf{\ourmethod{}} & \bfseries 78.4 & \bfseries 54.7 & 31 \\
      \bottomrule
    \end{tabular}
  \end{table}

  \vspace{0.8em}
  \begin{itemize}
    \setlength{\itemsep}{0.5em}
    \item 小目标子集的增益（\alert{+6.2}）显著高于总体（+3.7）
          \;$\Rightarrow$\; 提升确实来自小目标
    \item 速度从 \qty{33}{\fps} 降至 \qty{31}{\fps}，仍满足实时要求
  \end{itemize}

  \note{有人会问「为什么速度只掉 2 帧」——因为门控只在 3 个融合点上做，
        GAP 之后是 1×1 卷积，计算量可以忽略。}
\end{frame}

\begin{frame}{两个模块各自有效，且不相互抵消}
  \begin{columns}[T, onlytextwidth]
    \begin{column}{0.52\textwidth}
      \begin{table}
        \centering\small
        \begin{tabular}{cc S[table-format=2.1]}
          \toprule
          {门控} & {对齐} & {mAP / \%} \\
          \midrule
                     &            & 74.7 \\
          $\checkmark$ &            & 77.0 \\
                     & $\checkmark$ & 76.1 \\
          $\checkmark$ & $\checkmark$ & \bfseries 78.4 \\
          \bottomrule
        \end{tabular}
      \end{table}
    \end{column}
    \begin{column}{0.44\textwidth}
      \begin{itemize}
        \setlength{\itemsep}{0.7em}
        \item 门控单独贡献 \alert{+2.3}
        \item 对齐单独贡献 \alert{+1.4}
        \item 叠加得 +3.7，接近两者之和
      \end{itemize}
      \vspace{0.6em}
      {\footnotesize\color{AcadMuted}
       说明两个模块处理的是不同的误差来源。}
    \end{column}
  \end{columns}
\end{frame}

\begin{frame}{密集小目标场景的漏检明显减少}
  \begin{columns}[c, onlytextwidth]
    \begin{column}{0.48\textwidth}
      \centering
      \placeholder{\linewidth}{4.0cm}{figures/qual-baseline.pdf}\\[0.4em]
      {\footnotesize 基线：FPN + ATSS}
    \end{column}
    \begin{column}{0.48\textwidth}
      \centering
      \placeholder{\linewidth}{4.0cm}{figures/qual-ours.pdf}\\[0.4em]
      {\footnotesize \ourmethod{}（本文）}
    \end{column}
  \end{columns}
  \vspace{0.8em}
  \centering
  {\small 密集停放的车辆场景，红框为漏检。}
\end{frame}
